\section{Related Work}
\label{sec:related}

% Prose v1 (2026-06-10). Toutes les cles \cite ci-dessous existent et sont
% verifiees dans references.bib. Double-blind : OWNER cite en 3e personne.
% Cles encore manquantes (PDF non fourni), laissees en commentaire si besoin :
%   BERT (Devlin 2019), survey DL-NER (Li et al. 2022), InstructUIE, GPT-NER.

OPENER sits at the intersection of three lines of research: open-world and zero-shot named entity recognition (NER), pretrained entity representations refined by contrastive learning, and energy-aware machine learning.
We review each in turn and state how OPENER differs.

\subsection{Closed-set and Span-based NER}
\label{sec:related:closed}
Traditional NER is cast as supervised sequence labeling over a fixed inventory of entity types, typically with a pretrained encoder followed by a token- or span-level classifier.
Span-based formulations enumerate candidate spans and classify them directly, which simplifies the handling of nested and discontinuous entities \cite{nguyen2023span}.
These models reach high accuracy on in-domain benchmarks such as CoNLL-2003 \cite{tjongkimsang2003conll}, but their label set is frozen at training time.
Adding a new entity type requires re-annotating data and retraining the model, which is the central limitation that open-world NER seeks to remove.

\subsection{Open-world and Zero-shot NER}
\label{sec:related:openworld}
Open-world NER aims to recognise entity types that were not fixed at training time.
GLiNER \cite{zaratiana2024gliner} treats the candidate type names as queries to a bidirectional encoder, a DeBERTaV3 backbone \cite{he2023debertav3}, and scores every span against them in a single forward pass, offering a fast alternative to autoregressive tagging.
GNER \cite{ding2024gner} instead casts NER as text generation with a sequence-to-sequence model \cite{raffel2020t5} and shows that explicitly modeling negative, non-entity instances sharpens type boundaries.
Closest to our work, Genest et al. \cite{genest2025owner} propose OWNER, an unsupervised open-world pipeline that decouples mention detection from entity typing: it embeds detected mentions with an encoder refined by contrastive learning and then clusters them into dynamically named types.
Other approaches prompt a general-purpose large language model to perform extraction directly, as in ChatIE \cite{wei2024chatie}, or focus on recovering entities that detectors tend to miss \cite{cai2025handling}.
OPENER adopts the decoupled detect-then-type view of OWNER, but replaces the trained encoder and the unsupervised clustering with a pretrained, Matryoshka-truncatable embedding space and a balanced linear classifier, so that no encoder is trained from scratch.

\subsection{LLM-based and Few-shot NER}
\label{sec:related:llm}
A complementary line distils or instruction-tunes large language models for open NER.
UniversalNER \cite{zhou2024universalner} distils a chat model into a smaller LLaMA-based student \cite{touvron2023llama} over tens of thousands of entity types, while GoLLIE \cite{sainz2024gollie} conditions a code language model on annotation guidelines to follow unseen schemas.
Few-shot methods reduce supervision further through prompt optimization \cite{tang2024fsponer}, knowledge-graph-guided contrastive learning \cite{zhang2024kcl}, or token-level contrastive objectives such as CONTaiNER \cite{das2022container}.
These methods are accurate but rest on models with billions of parameters, whose inference cost, latency and energy use are substantial, and which may not even fit on commodity hardware.
This cost is precisely what motivates the lightweight design and the energy-aware evaluation of OPENER.

\subsection{Entity Representations and Contrastive Learning}
\label{sec:related:repr}
OPENER relies on a pretrained sentence embedder rather than a task-specific encoder.
Nomic Embed \cite{nussbaum2024nomic} provides an open, long-context model trained with the Matryoshka objective \cite{kusupati2022matryoshka}, whose representations can be truncated to lower dimensions without retraining, trading quality for compute.
Such embedders build on the Sentence-BERT line of work \cite{reimers2019sbert}, and we sharpen the entity-type geometry with a triplet margin objective in the spirit of FaceNet \cite{schroff2015facenet}.
Embedding labels and entities in a shared space \cite{wang2018joint}, modeling entity types with Gaussian components \cite{zhang2025contrastivegaussian}, generating synthetic entity-bearing text for augmentation \cite{hu2023entity}, and using mixtures to augment training data \cite{abbahaddou2024gmmgnn} are recurring ideas in low-resource NER, all grounded in classical Gaussian mixture models fitted by expectation-maximization \cite{reynolds2009gmm, dempster1977em}.
Unlike methods that train a dedicated encoder, OPENER keeps the pretrained space largely frozen and only applies a light contrastive fine-tuning, which tests whether a general-purpose embedding is already discriminative enough for entity typing.

\subsection{Efficient and Green AI}
\label{sec:related:green}
A growing body of work argues that efficiency should be reported alongside accuracy.
Strubell et al. \cite{strubell2019energy} quantify the financial and environmental cost of training NLP models, Schwartz et al. \cite{schwartz2020greenai} advocate efficiency as a first-class evaluation criterion, and Patterson et al. \cite{patterson2022carbon} analyse the carbon footprint of machine-learning training and how to reduce it.
Tools such as CodeCarbon \cite{lacoste2019codecarbon} make per-run energy and carbon estimates practical.
Open-world NER systems, however, are still compared almost exclusively on quality.
OPENER instead reports a three-axis comparison of quality, measured by Adjusted Mutual Information \cite{vinh2010ami}, inference speed, and energy, placing the frugality of each method on an equal footing with its accuracy.

In summary, existing open-world NER methods either train a dedicated encoder \cite{genest2025owner} or query costly large language models \cite{zhou2024universalner, sainz2024gollie}, and they are compared on accuracy alone.
OPENER fills both gaps: it reuses pretrained components with only a light contrastive step and a balanced linear classifier, and it is evaluated jointly on quality, speed and energy.
